\documentclass[11pt,a4paper]{article}
\usepackage[margin=2.5cm]{geometry}
\usepackage{longtable,array,booktabs,enumitem,hyperref}
\setlength{\parindent}{0pt}
\setlength{\parskip}{6pt}

\title{VeriPilot: Work Split and Plan to 20 November}
\author{Team of five}
\date{Draft, 3 Nov 2026}

\begin{document}
\maketitle

\section{Goal}
Show a working demo by \textbf{20 November 2026}. The demo runs one command on a small design and shows lint results, per-seed test results, coverage gaps and a health score on a dashboard. The guiding rule stays the same: \emph{AI suggests, the simulator verifies.}

Step 1 is a thin end-to-end skeleton on a single design (a FIFO). Each stage can be crude at first. All five of us work in parallel because each role owns different directories.

\section{Related work}
Two public repositories share part of our name or scope. Neither covers our full loop.
\begin{itemize}[nosep]
  \item \textbf{YihanWn/VeriPilot}: a research prototype that localises and repairs Verilog bugs with an LLM and a Python golden model. It has no lint, coverage, regression or dashboard. It has one commit and no licence. It shares our name.
  \item \textbf{fangzhigang32/RTLBench}: an ICCD 2025 benchmark of 160 designs that scores LLM-generated RTL (syntax, function, lint, readability). It measures code generation, not verification of a user's design. MIT licensed.
\end{itemize}
We differ by connecting lint, safe fixes, regression, coverage guidance, failure triage and progress tracking in one loop. Because of the name clash, the name decision in Section~\ref{sec:decisions} is urgent.

\section{Roles}
\begin{longtable}{p{2.6cm} p{3.2cm} p{9.2cm}}
\toprule
\textbf{Role} & \textbf{Owns} & \textbf{Responsibilities} \\
\midrule
\endhead
Hardware lead & \texttt{benchmarks/} & Choose 4 to 6 small open designs (FIFO, UART, SPI, timer) and get them running on Verilator. Write testbenches. Plant bugs with a ground-truth file per design. Later, check that coverage and bug-detection numbers are correct. \\
Software 1 (fixer) & \texttt{lint/}, \texttt{autofix/} & Run the linter and parse its output. Apply a fix only if the tests still pass, otherwise revert. Add a small local LLM to draft fixes once the rule-based version works. \\
Software 2 (tester) & \texttt{regress/}, \texttt{coverage\_advisor/}, \texttt{triage/} & Run N seeds and record pass/fail per seed. Parse the coverage report and ask the LLM for stimulus that reaches uncovered cases, keeping it only if coverage rises. Build the flaky-versus-repeatable check. \\
Product & \texttt{dashboard/}, user side of \texttt{docs/} & Interview 3 to 5 target users. Define and build the dashboard and the single health score from the results file. Write the user README and design the demo flow. \\
Project/Business & repo, process, submission & Close the blocking decisions (name, licence, repo location, contribution rules). Run the weekly meeting and mentor contact. Own the results JSON schema. Run the baseline-versus-VeriPilot comparison. Build the demo run-sheet and slides. Submit. \\
\bottomrule
\end{longtable}

Software 2 carries the heaviest load. If it becomes too much, move \texttt{triage/} to Software 1 once autofix is stable. Hardware has a lighter middle stretch and can help test the other stages.

\section{How the roles connect}
The roles connect through one shared \textbf{results JSON} written by each stage: lint count, per-seed results, uncovered bins and the bug list. We agree its format in the first two days (all five, about one hour). After that nobody blocks anyone.

Two dependencies to watch:
\begin{itemize}[nosep]
  \item Software 1 and 2 need a working design from Hardware. The first FIFO is due within about three days. Until then they can use any public example.
  \item Project/Business needs Hardware's planted-bug ground truth to score results.
\end{itemize}

\section{Timeline}
\begin{longtable}{p{1.9cm} p{2.5cm} p{2.5cm} p{2.5cm} p{2.5cm} p{2.5cm}}
\toprule
\textbf{Dates} & \textbf{Hardware} & \textbf{Software 1} & \textbf{Software 2} & \textbf{Product} & \textbf{Project/ Business} \\
\midrule
\endhead
3--8 Nov & FIFO running, bugs planted & Lint parser, rule-based safe-fix loop & Seed runner, coverage parser & User interviews, dashboard sketch & Schema, decisions, repo setup \\
9--14 Nov & Second design (UART) & LLM drafts fixes & Coverage advisor, triage & Working dashboard & Baseline run \\
15--18 Nov & Verify numbers & Fix bugs, tune & Fix bugs, tune & Polish, demo flow & Full comparison, slides \\
19--20 Nov & \multicolumn{5}{l}{Rehearse and submit together} \\
\bottomrule
\end{longtable}

\section{Definition of done for step 1}
Running \texttt{veripilot run benchmarks/fifo} produces a lint count, pass/fail per seed, a list of uncovered bins and a single results JSON file.

\section{Decisions needed this week}
\label{sec:decisions}
\begin{itemize}[nosep]
  \item Who holds which role, and the owner of each directory.
  \item The results JSON schema.
  \item Project name (search GitHub and PyPI; fallbacks: BugBloom, CoverQuest, SeedSense), licence and repo location.
  \item Testbench style (plain SystemVerilog or cocotb) and Python version.
  \item Meeting day and communication channel.
  \item Which small local model to use for fixes (can wait until 9 Nov).
\end{itemize}

\section{Known challenges}
\begin{itemize}[nosep]
  \item \textbf{Slow feedback.} A hardware regression is slower than a software unit test, so each round gets a small budget of LLM attempts.
  \item \textbf{Weak oracle.} A test that passes and adds coverage may still check nothing. We rely on self-checking testbenches or a reference model, and on planted bugs to confirm the tests catch real faults.
  \item \textbf{Coverage depth.} Verilator provides line and toggle coverage. Functional coverage bins need extra work.
\end{itemize}

\end{document}
